\begin{table}[H]
\centering
\footnotesize
\caption{Categories of Respiratory Sounds - adopted from \cite{pasterkamp_respiratory_1997}}
\label{tab:resp_sounds}

\begin{tabularx}{\textwidth}{lXXXX}
\toprule
\textbf{Respiratory Sound} & \textbf{Mechanisms} & \textbf{Origin} & \textbf{Acoustics} & \textbf{Relevance} \\
\midrule

\multicolumn{5}{l}{\textbf{Basic sounds}} \\
\addlinespace

Normal lung sound &
Turbulent flow vortices, unknown mechanisms &
Central airways (expiration), lobar to segmental airways (inspiration) &
Low-pass filtered noise (range $<100$ to $>1000$ Hz) &
Regional ventilation, airway caliber \\
\addlinespace

Normal tracheal sound &
Turbulent flow, flow impinging on airway walls &
Pharynx, larynx, trachea, large airways &
Noise with resonances (range $<100$ to $>3000$ Hz) &
Upper airway configuration \\
\midrule

\multicolumn{5}{l}{\textbf{Adventitious sounds}} \\
\addlinespace

Wheeze &
Airway wall flutter, vortex shedding &
Central and lower airways &
Sinusoid (range $\sim100$ to $>1000$ Hz; duration typically $>80$ ms) &
Airway obstruction, flow limitation \\
\addlinespace

Rhonchus &
Rupture of fluid films, airway wall vibrations &
Larger airways &
Series of rapidly dampened sinusoids (typically $<300$ Hz; duration $>100$ ms) &
Secretions, abnormal airway collapsibility \\
\addlinespace

Crackle &
Airway wall stress-relaxation &
Central and lower airways &
Rapidly damped wave deflection (duration typically $<20$ ms) &
Airway closure, secretions \\
\bottomrule
\end{tabularx}

\vspace{0.3cm}
\footnotesize{
* This table lists only the major categories of respiratory sounds and does not include other sound such as squawks, friction rubs, grunting, snoring, or cough. Current concepts on sound mechanisms and origin may be incomplete and unconfirmed
}
\end{table}